\chapter{Distributed Electronics and Event Processing}

Each band has a specific antenna, preselector, limiter, LNA, converter when needed, coherent ADC, FPGA/ASIC and circular memory. The local chain is:
\[
\text{antenna}\rightarrow\text{RF front-end}\rightarrow\text{ADC}\rightarrow\text{DDC/PFB/FFT}\rightarrow\text{detector}\rightarrow\text{event}.
\]
The system performs three reductions: band-level detection, face-level fusion and global radome fusion. Raw samples remain local and are preserved around triggers. Normal network traffic consists of event descriptors, telemetry, summaries and selected I/Q captures rather than permanent aggregate ADC streaming.

For C3 storage screening, the aviation paths use a proposed 8 MS/s complex capture with 16-bit I and 16-bit Q for each of the simultaneous UAT and 1090ES channels (\texttt{RF-012}). A selected direct UHF waveform uses a proposed 25 MS/s complex channel (\texttt{RF-013}), while the bistatic protocol uses two coherent 25 MS/s channels for reference and surveillance (\texttt{RF-014}). These rates produce 512 Mbit/s, 800 Mbit/s and 1.6 Gbit/s per node respectively before framing or compression. Ten-second raw event windows therefore require 0.64 GB, 1.00 GB and 2.00 GB per node. They are capacity baselines, not claims that each waveform occupies the full sampled bandwidth.

Noise figure, IP3 and usable dynamic range remain deliberately unassigned (\texttt{RF-015}) until a measured site-RFI distribution, antenna gain pattern, preselector rejection and cascaded component budget exist. Component selection must close sensitivity and blocker survival together; an isolated low-noise figure is not sufficient.

Every antenna port assigned to a band has its own shielded acquisition module containing at least one ADC and one band ASIC. A dual-polarized aperture therefore uses two independent enclosures, one for each orthogonal coherent channel (\texttt{MEC-004}); sharing the sampling reference does not mean sharing an uncontrolled RF/digital cavity. The modules are separated by conductive partitions, RF gaskets and controlled cable/fibre feedthroughs. A central face module, the FFASIC, receives events from all band modules, performs intra-face correlation and interfaces with timing, thermal monitoring, power and the optical link.

Thermally, each shielded module is mounted to a conductive cold plate connected to a controlled heat path. The RF cavity, digital processing electronics and power conversion are physically partitioned so that heat sources and digital EMI do not share an uncontrolled volume with the sensitive RF aperture. Temperature sensors are attached to the LNA, ADC, ASIC, clock and cold plate; thermal drift is recorded as metadata and triggers recalibration when the validated range is exceeded.

\begin{figure}[H]
\centering
\includegraphics[width=0.95\textwidth]{fig03_modulo_face.png}
\caption{Triangular RF face module and exploded stack: the left triangle represents aperture, calibration and structural paths; the right coloured plates represent successive RF, dielectric, shielding, processing, thermal and service layers.}
\label{fig:fig03_modulo_face}
\end{figure}

\begin{figure}[H]
\centering
\includegraphics[width=0.95\textwidth]{fig06_cadeia_rf.png}
\caption{RF and edge-processing chain. The upper sequence represents antenna-to-event processing; the three lower blocks and vertical arrows represent calibration injection, common timing and supervisory telemetry.}
\label{fig:fig06_cadeia_rf}
\end{figure}
